\section{An interactive review of the literature}
\label{sec:literature}

Choosing hyper-parameters, datasets and baselines for a Tsetlin machine experiment means
reading the literature for numbers it rarely tabulates in one place. The package therefore
ships a catalog of every dataset and every reported number from \litNpapers\ Tsetlin-machine
papers published between \litYearMin\ and \litYearMax. The papers span the foundational
variants, image classification, segmentation, medical and clinical applications, hardware
accelerators, generative and sequence models, and privacy-preserving learning. The catalog
covers both the methods each paper proposed and the baselines it compared against.

\begin{center}
\small\sffamily\sloppy
\captionof{table}{What the literature catalog records.}
\label{tab:literature}
\begin{tabularx}{\linewidth}{@{}L{2.5cm}R{1.4cm}>{\raggedright\arraybackslash}X@{}}
\ttoprule
\thead{Table} & \thead{Rows} & \thead{One row per\ldots}\\
\ttmidrule
\code{papers} & \litNpapers & paper: venue, arXiv/DOI, code, the Tsetlin-machine variants tested,
  the feature it introduces, and a caution flag where a reported number has a specific
  reliability problem (\litNflagged\ papers)\\
\code{datasets} & \litNdatasets & dataset: data type, task, classes, input shape, download link
  (\litNdownload\ public, each one checked) and the variants that used it\\
\code{dataset\_usages} & \litNusages & (paper, dataset) pair: the train\slash val\slash test sizes, the
  split exactly as reported, and how the inputs were Booleanized\\
\code{results} & \litNresults & reported number: model, metric (\litNmetrics\ distinct), value,
  unit, split, configuration, and the table or figure it came from; \litNtm\ for
  Tsetlin-machine models, \litNbaseline\ for baselines\\
\ttbotrule
\end{tabularx}
\end{center}

\subsection{How it was built}

Each paper has one JSON record in \code{benchmarks/literature/papers/}, and those records are
the source of truth. Numbers were transcribed from the paper's tables and text. Where an
extracted table was garbled, they were checked against the rendered PDF page. Values that
appear only in plots were left out rather than estimated. Dataset names are unified to one
identifier (\code{kmnist} and \code{kuzushiji-mnist} are the same dataset), and a field a paper
does not state stays empty. \code{build\_catalog.py} turns the records into the SQLite file
\code{docs/assets/literature/literature.db} (\litDbKiB~KiB) and a flat CSV.

\subsection{Using it}

The documentation's \emph{Literature results} page opens the database in the browser with
\code{sql.js}, so filtering, grouping and ranking run as SQL on the reader's machine, with no
server. Its tabs are the results themselves, a best-per-dataset leaderboard, the datasets with
their download links and per-paper splits, the papers, and a SQL console. Within any grouping
the page stars the best value. Values are compared only when they share the metric, unit and
split, and a metric's direction decides what counts as best: higher for accuracy, $F_1$ or
throughput, lower for error, latency, energy or area. The same database backs a Hugging Face
Space\footnote{\url{https://huggingface.co/spaces/vlbthambawita/TsetlinMachines}}, and the
file can be opened with any SQLite client:

\begin{lstlisting}[style=ttshell,basicstyle=\ttfamily\scriptsize\color{ttink}]
SELECT p.year, p.title, r.model, MAX(r.value) AS accuracy
FROM results r JOIN papers p USING (paper_id)
WHERE r.dataset_id = 'cifar-10' AND r.metric = 'accuracy' AND r.is_tm = 1
GROUP BY r.paper_id ORDER BY accuracy DESC;
\end{lstlisting}

\begin{figure}[htbp]
\centering
\resizebox{\linewidth}{!}{% Figure: the literature catalog. Data: data/lit-top.csv and data/lit-best.csv, both written by
% data/regen_literature.py from docs/assets/literature/literature.db.
\pgfplotstableread[col sep=comma]{data/lit-top.csv}\littop
\pgfplotstableread[col sep=comma]{data/lit-best.csv}\litbest
\begin{tikzpicture}
\begin{axis}[
    ttaxis,
    name=top,
    width=0.47\linewidth, height=6.4cm,
    xbar, bar width=6pt,
    xmin=0,
    xlabel={papers using the dataset},
    ytick=data, y dir=reverse,
    yticklabels from table={\littop}{dataset},
    yticklabel style={font=\ttfamily\scriptsize, color=ttink},
    ymajorgrids=false,
    enlarge y limits=0.06,
    nodes near coords,
    nodes near coords style={font=\sffamily\tiny, color=ttdeep},
    title={Most-used datasets},
  ]
  \addplot[ttbar] table[x=papers, y expr=\coordindex] {\littop};
\end{axis}
\begin{axis}[
    ttaxis,
    at={($(top.east)+(1.9cm,0)$)}, anchor=west,
    width=0.53\linewidth, height=6.4cm,
    ybar, bar width=6.5pt,
    ymin=75, ymax=101,
    ylabel={best reported test accuracy (\%)},
    xtick=data,
    xticklabels from table={\litbest}{dataset},
    xticklabel style={font=\sffamily\scriptsize, color=ttink, rotate=35, anchor=north east},
    xmajorgrids=false,
    enlarge x limits=0.1,
    legend style={at={(0.5,1.02)}, anchor=south, legend columns=2},
    title={TM vs.\ the baselines TM papers report},
    title style={yshift=13pt},
  ]
  \addplot[ttbar] table[x expr=\coordindex, y=tm] {\litbest};
  \addlegendentry{best TM model}
  \addplot[ttbarcpu] table[x expr=\coordindex, y=baseline] {\litbest};
  \addlegendentry{best non-TM baseline}
\end{axis}
\end{tikzpicture}
}
\caption{The literature catalog. \emph{Left:} the datasets used by the most papers.
  \emph{Right:} the best test accuracy reported for a Tsetlin-machine model next to the best
  non-TM baseline reported in the same body of papers. Flagged papers and federated
  (per-client) accuracies are excluded. The baselines are those TM papers chose to compare
  against, not the state of the art in each field.}
\label{fig:literature}
\end{figure}

Figure~\ref{fig:literature} shows what the catalog makes easy to see. Tsetlin machines are
ahead of the baselines they are reported against on Fashion-MNIST and Reuters~R8, close on
MNIST and Kuzushiji-MNIST, and behind on CIFAR-10 and IMDb.

\begin{ttcallout}[Reading the numbers]
Every value is \emph{as reported by its authors}; none was reproduced for the catalog. Papers
differ in splits, Booleanization, clause budgets and epochs. Some report the peak test
accuracy over training, and results measured before 0.3.0 made the clause budget a hard cap
are not directly comparable with \ttlibname runs. Every row carries its configuration, source
location and notes, and the page shows them on click.
\end{ttcallout}
